
\clearpage
\phantomsection
\addcontentsline{toc}{chapter}{CHƯƠNG 1: TỔNG QUAN VÀ QUY TRÌNH THIẾT KẾ}
{\LARGE\bfseries CHƯƠNG 1: TỔNG QUAN VÀ QUY TRÌNH THIẾT KẾ}\par\vspace{1em}

\phantomsection
\addcontentsline{toc}{section}{1.1. Giới thiệu hệ thống FamilyConnect}
{\Large\bfseries 1.1. Giới thiệu hệ thống FamilyConnect}\par\vspace{0.7em}

FamilyConnect phục vụ các dòng họ/gia đình quản lý: tài khoản và phân quyền; cây gia phả (chi/nhánh, quan hệ cha mẹ – con, hôn nhân); không gian cộng đồng (bài đăng, bình luận, cảm xúc, media); sự kiện gia đình (họp mặt, giỗ, sinh nhật kèm RSVP); danh bạ nghề nghiệp/học vấn của thành viên; di sản số (tài liệu lịch sử, câu chuyện, thành viên tiêu biểu, album ảnh); dịch vụ AI (tìm kiếm ngữ nghĩa, trợ lý tri thức, giải thích quan hệ họ hàng); thống kê báo cáo; và quản trị hệ thống (nhật ký, cấu hình, sao lưu, kiểm duyệt).

\phantomsection
\addcontentsline{toc}{section}{1.2. Áp dụng chu kỳ sống của CSDL}
{\Large\bfseries 1.2. Áp dụng chu kỳ sống của CSDL}\par\vspace{0.7em}

Việc thiết kế không bắt đầu bằng câu hỏi "tạo bảng nào" mà đi qua đầy đủ các giai đoạn: Phân tích yêu cầu → Thiết kế quan niệm → Thiết kế logic → Thiết kế vật lý → Triển khai → Vận hành. Bảng dưới đây áp dụng dòng chảy input–output của giáo trình vào đề tài FamilyConnect.

\begin{longtable}{|p{0.3\textwidth}|p{0.3\textwidth}|p{0.3\textwidth}|}
\hline
\textbf{Giai đoạn} & \textbf{Input áp dụng cho FamilyConnect} & \textbf{Output áp dụng cho FamilyConnect} \\ \hline
\endfirsthead
\hline
\textbf{Giai đoạn} & \textbf{Input áp dụng cho FamilyConnect} & \textbf{Output áp dụng cho FamilyConnect} \\ \hline
\endhead
Phân tích yêu cầu & Đặc tả 9 nhóm chức năng, quy tắc nghiệp vụ (VD: một thành viên có thể chưa có tài khoản) & Từ điển thuật ngữ nghiệp vụ, danh sách ràng buộc \\ \hline
Thiết kế quan niệm & Từ điển nghiệp vụ ở trên & ERD/EER, bản số, thực thể yếu \\ \hline
Thiết kế logic & ERD đã xác nhận & Lược đồ quan hệ, PK/FK, chứng minh chuẩn hóa 3NF/BCNF \\ \hline
Thiết kế vật lý & Lược đồ quan hệ + SQL Server & Script T-SQL, kiểu dữ liệu, chỉ mục, ràng buộc  \\ \hline
Triển khai & Script T-SQL & CSDL FamilyConnectDB chạy trên container Docker SQL Server 2022 \\ \hline
\end{longtable}


\phantomsection
\addcontentsline{toc}{section}{1.3. Môi trường triển khai}
{\Large\bfseries 1.3. Môi trường triển khai}\par\vspace{0.7em}

Hệ quản trị CSDL được chọn là Microsoft SQL Server 2022, cài đặt qua Docker để bảo đảm môi trường đồng nhất:

\begin{lstlisting}[language=bash]
docker pull mcr.microsoft.com/mssql/server:2022-latest
docker run --name sqlserver2022 --hostname sqlserver2022 \
  -e "ACCEPT_EULA=Y" -e "MSSQL_SA_PASSWORD=FamilyConnect@2026" \
  -p 1433:1433 -v familyconnect_data:/var/opt/mssql \
  -d mcr.microsoft.com/mssql/server:2022-latest
\end{lstlisting}

Sau khi container chạy, script schema\_sqlserver.sql (Phụ lục A) được nạp qua sqlcmd hoặc Azure Data Studio để tạo CSDL FamilyConnectDB.